\chapter*{Conclusion}
\addcontentsline{toc}{chapter}{Conclusion}

This thesis proposes and evaluates a data-driven dynamic pricing model for short-term accommodation providers. The model operates in two stages: a demand forecasting stage based on SARIMA, which predicts future occupancy levels from historical reservation data, and a pricing stage that derives price recommendations as quantiles of the conditional price distribution estimated via kernel density estimation of historical occupancy–price pairs. The design deliberately avoids reliance on external market data or explicit estimation of price elasticity, making the approach accessible to small accommodation operators whose primary data asset is their own booking history.

The model was applied to an anonymized dataset of 3{,}787 reservation records covering 13 apartments in Prague, Czech Republic, over a period of 394 days. The $\mathrm{SARIMA}(2,1,1)(1,0,2)_7$ model was selected automatically using AIC and achieved a Mean Absolute Error of 7.29 percentage points on the 30-day hold-out test period, representing a 22.1\% improvement over a seasonal naive baseline. Analysis of forecast errors revealed a systematic smoothing: the model compresses the range of predicted occupancy, over-predicting on low-demand days and under-predicting on high-demand days, which directly propagates into the pricing stage. Evaluation of price recommendations at the median quantile $\alpha=0.5$ across all 13 apartments and the full test period showed a mean signed error of only $-4.2\%$. However, MAE of $28.4\%$ and a standard deviation of $34.7\%$ indicate significant variability in pricing accuracy across different days and apartments. The results confirm that the proposed model meets its objective of providing an accessible, data-driven pricing support tool that requires no external market data, and the identified limitations offer a roadmap for extending its practical applicability.

Several directions for future work emerge from this study. Incorporating external variables such as public holidays, local events, or booking lead times into the demand forecasting component could reduce systematic forecasting errors during demand spikes. Allowing the quantile parameter $\alpha$ to vary dynamically as a function of the forecasting horizon or predicted demand uncertainty would enable more flexible and context-sensitive pricing strategies. Ensuring precise underlying data remains a critical part of model performance.
